%Set paragraphs
\setlength{\parskip}{0.25em} %gap between paragraph
\setlength{\parindent}{2.1em} %indentation of paragraph

%\usepackage[switch, modulo]{lineno}
\usepackage[]{lineno}
%\usepackage[modulo]{lineno}

\usepackage{amsmath}
\usepackage{bbm}
\usepackage{amsfonts} 
%additional math fonts
\usepackage[mathscr]{euscript} %additional math
%for expectation symbol
%\DeclareMathOperator*{\E}{\mathbb{E}}

\usepackage{graphicx}
\usepackage{verbatim}%for multiline comments
\usepackage[dvipsnames,svgnames,x11names,hyperref]{xcolor}
\usepackage{tabularx} %%get textwidth tables
\usepackage{longtable}
\usepackage{siunitx}
\newcolumntype{d}{S[input-symbols = ()]}

% for landscape
\usepackage{lscape}

\usepackage{booktabs} % for using midrule in table
\usepackage{xspace} % To get the spacing after macros right
%------------------------------------------------
\usepackage{makecell} % To get multiline table cells
%------------------------------------------------
\usepackage{array,multirow,graphicx} % to get rotating text in table
\usepackage{rotating} %for rotated figs
%------------------------------------------------
\usepackage{mparhack} % To get marginpar right
\frenchspacing % Reduces space after periods to make text more compact
\usepackage{varioref} % More descriptive referencing

%%%shortcut for graphics location
\graphicspath{ {images/} }

%For double spacing
\usepackage{setspace}
\onehalfspacing


\usepackage[T1]{fontenc}
\usepackage[british]{babel}



%margins
\usepackage[a4paper, left= 1 in , right= 1 in, top= 1 in, bottom= 1 in]{geometry}
%\usepackage[left= 1 in , right= 1 in, top= 1 in, bottom= 1 in]{geometry}


\usepackage[
	backend = biber,
	bibencoding = latin1,
	style=apa]{biblatex}
\DeclareLanguageMapping{british}{british-apa}
\addbibresource{backmatter/RIFdecomposition.bib}

\usepackage{fancyhdr}


\usepackage{footnotebackref}
\usepackage{epsfig,graphics}
\usepackage{etoolbox}
%\overfullrule=5pt %%%%%%Remove this to remove badbox indicator

\usepackage{appendix}

%making the list of symbols 
\usepackage[]{nomencl}
\makenomenclature

\nomlabelwidth= 4em % to change the gap between the symbol and text
\renewcommand{\nomname}{List of Symbols}
\usepackage{etoolbox}


%%%%Code formatting 
%\usepackage{minted}
\usepackage{csquotes}
\usepackage{lmodern}

%%%%%%%%%%tikz

% Include TikZ and PGF packages for high-quality graphics, schematics
% and plots. This is optional; at the current time, when run with
% latex to create a .dvi file, the xdvi viewer will produce incorrect
% formatting for the TikZ figures.  If the .dvi file is converted to
% pdf using "dvipdf" the resulting pdf file is correct.  If this
% example is used with  pdflatex, the resulting TikZ figures in the
% output look fine.
\usepackage{tikz}       		% The base tikz+pgf package
\usetikzlibrary{arrows,shapes, 
            chains, positioning}  %Tikz extension 
\usepackage{pgfplots}			% Tikz-based package for making plots 
\pgfplotsset{compat=1.6}      % This *might* be necessary for your
                               % version of pgfplots.


%Using hyperref at the last.


\usepackage{hyperref}
\hypersetup{
%draft, % Uncomment to remove all links (useful for printing in black and white)
linktocpage,
colorlinks=true, breaklinks=true, %bookmarks=true, bookmarksnumbered,
urlcolor=[RGB]{128,0,0}, linkcolor=[RGB]{25,25,112}, citecolor=[RGB]{72,61,139},  % Link colors
pdftitle={RIF Decomposition}, % PDF title
pdfauthor={Sandeep, Resham, Sanjeet and Hemanta}, % PDF Author
pdfsubject={Labor economics}, % PDF Subject
pdfkeywords={Informal wage gap, RIF Decomposition}, % PDF Keywords
pdfcreator={Sandeep, Resham and Sanjeet}, % PDF Creator
}

\makeatletter
\def\@footnotecolor{red}
\define@key{Hyp}{footnotecolor}{%
 \HyColor@HyperrefColor{#1}\@footnotecolor%
}
\patchcmd{\@footnotemark}{\hyper@linkstart{link}}{\hyper@linkstart{footnote}}{}{}
\makeatother
\hypersetup{footnotecolor=[RGB]{255,69,0}}


